\begin{theorem}[秩一耦合参数族的全实轴惯性相图与唯一零模穿越]\label{thm:xi-replica-softcore-exceptional-rankone-coupling-full-axis-inertia}
沿用定理 \ref{thm:xi-replica-softcore-exceptional-rankone-coupling-inertia-rigidity} 的记号，并把参数域扩展到全部实数 \(t\in\RR\)：
\[
A^{(q)}(t):=\frac12\Bigl(C^{(q)}+t\,\mathbf 1(b^{(q)})^{\mathsf T}\Bigr).
\]
记 \(\mathsf{In}(M):=(n_+(M),n_-(M),n_0(M))\) 为实矩阵 \(M\) 的正、负、零特征值个数组成的惯性三元组。
则对任意实数 \(t\) 都有
\[
\boxed{
\mathsf{In}\bigl(A^{(q)}(t)\bigr)
=
\begin{cases}
\left(\left\lceil\frac{q+1}{2}\right\rceil,\ \left\lfloor\frac{q+1}{2}\right\rfloor,\ 0\right),& t>-1,\\[6pt]
\left(\left\lceil\frac{q+1}{2}\right\rceil-1,\ \left\lfloor\frac{q+1}{2}\right\rfloor,\ 1\right),& t=-1,\\[6pt]
\left(\left\lceil\frac{q+1}{2}\right\rceil-1,\ \left\lfloor\frac{q+1}{2}\right\rfloor+1,\ 0\right),& t<-1.
\end{cases}}
\]
换言之，参数族 \(A^{(q)}(t)\) 在全实轴上仅发生一次且恰为一维的零模穿越，穿越点唯一位于 \(t=-1\)，并且穿越只把一个正特征值翻转为负特征值。
\end{theorem}
\begin{proof}
沿结论 \ref{con:xi-symq-golden-involution-reciprocity} 记
\[
W_q:=\mathrm{diag}\Bigl(\binom{q}{0}^{-1},\binom{q}{1}^{-1},\dots,\binom{q}{q}^{-1}\Bigr),
\qquad
u_q:=W_q^{-1/2}\mathbf 1.
\]
由于 \(b^{(q)}=W_q^{-1}\mathbf 1\)，有
\[
W_q^{-1/2}\,\mathbf 1(b^{(q)})^{\mathsf T}\,W_q^{1/2}
=
u_qu_q^{\mathsf T}.
\]
另一方面，由结论 \ref{con:xi-symq-golden-involution-reciprocity} 中同一矩阵的加权自伴随关系，
\[
C^{(q)\mathsf T}W_q=W_qC^{(q)},
\]
可知
\[
S_q:=W_q^{-1/2}C^{(q)}W_q^{1/2}
\]
为实对称矩阵。
因此
\[
\widetilde A^{(q)}(t)
:=
W_q^{-1/2}A^{(q)}(t)W_q^{1/2}
=
\frac12\Bigl(S_q+t\,u_qu_q^{\mathsf T}\Bigr)
\]
是一个与 \(A^{(q)}(t)\) 相似的实对称矩阵。
故 \(A^{(q)}(t)\) 具有全实谱，且其惯性可由 \(\widetilde A^{(q)}(t)\) 计算。

定理 \ref{thm:xi-replica-softcore-exceptional-rankone-coupling-inertia-rigidity} 中对行列式的推导是关于 \(t\) 的代数恒等式，因而同样适用于全部实数 \(t\)。
于是
\[
\det A^{(q)}(t)=2^{-(q+1)}\det(C^{(q)})(1+t).
\]
因此唯一可能的奇异参数是 \(t=-1\)，并且该零点为一次零点。
于是对称矩阵族 \(\widetilde A^{(q)}(t)\) 在两个连通区间 \((-\infty,-1)\) 与 \((-1,\infty)\) 上惯性各自保持常值。

在 \(t=0\) 时，
\[
A^{(q)}(0)=\frac12C^{(q)},
\qquad
\mathrm{spec}\bigl(A^{(q)}(0)\bigr)
=
\left\{\frac{(-1)^j\varphi^{q-2j}}{2}\right\}_{j=0}^{q},
\]
故
\[
\mathsf{In}\bigl(A^{(q)}(0)\bigr)
=
\left(\left\lceil\frac{q+1}{2}\right\rceil,\ \left\lfloor\frac{q+1}{2}\right\rfloor,\ 0\right).
\]
从而对全部 \(t>-1\) 都有同一惯性三元组。

接着考察 \(t=-1\)。
由于
\[
\widetilde A^{(q)}(-1)
=
\widetilde A^{(q)}(0)-\frac12u_qu_q^{\mathsf T},
\]
这是可逆矩阵 \(\widetilde A^{(q)}(0)\) 的秩一扰动，故其零空间维数至多为 \(1\)。
另一方面 \(\det \widetilde A^{(q)}(-1)=\det A^{(q)}(-1)=0\)，故
\[
n_0\bigl(A^{(q)}(-1)\bigr)=1.
\]
又因为 \(t>-1\) 一侧的惯性已知，遂得
\[
\mathsf{In}\bigl(A^{(q)}(-1)\bigr)
=
\left(\left\lceil\frac{q+1}{2}\right\rceil-1,\ \left\lfloor\frac{q+1}{2}\right\rfloor,\ 1\right).
\]

最后确定穿越方向。
若 \(t_2>t_1\)，则
\[
\widetilde A^{(q)}(t_2)-\widetilde A^{(q)}(t_1)
=
\frac{t_2-t_1}{2}\,u_qu_q^{\mathsf T}\succeq 0.
\]
由 Courant--Fischer 极小极大原理，\(\widetilde A^{(q)}(t)\) 的每个特征值都是 \(t\) 的非减函数。
因此当 \(t\) 从 \(-1+\varepsilon\) 向左穿越到 \(-1-\varepsilon\) 时，只能有一个特征值由正侧下降穿过 \(0\)；不可能有负特征值向上穿过 \(0\)。
再结合 \(t=-1\) 处零空间恰为一维，即得对全部 \(t<-1\)
\[
\mathsf{In}\bigl(A^{(q)}(t)\bigr)
=
\left(\left\lceil\frac{q+1}{2}\right\rceil-1,\ \left\lfloor\frac{q+1}{2}\right\rfloor+1,\ 0\right).
\]
证毕。
\end{proof}
